\section{Discussion}
\label{sec:discussion}

\para{Answer to the hypothesis.} For the two-stage review$\to$revise loop the answer is a
clear \emph{no}: it does not converge to a fixed point within a practical horizon; it enters a
non-converging \emph{expansion regime}. The hypothesis \emph{is} essentially true for the
single-stage variant---self-refine relaxes to a soft fixed point---which is the crucial
contrast: the explicit, regenerated review is what prevents convergence, not iteration
per se.

\para{Why the review stage matters.} A ``make this better'' editor can declare a document
locally optimal and stop editing ($\dt\to 0$). A ``peer reviewer'' role is adversarial by
construction: it is prompted to \emph{find weaknesses}, and a capable model can always find
some. Each critique is fresh edit pressure, and because the cheapest way to ``address'' most
critiques is to add explanatory material, the operator $F$ behaves less like a contraction
toward a fixed point and more like a \emph{growth map}. This reconciles the two camps in the
literature: the convergent relaxation of~\citet{wu2026converge} holds for single-stage
self-refine, while the skeptics' warnings---self-bias amplification~\citep{xu2024pride},
feedback friction~\citep{feedbackfriction2025}, recursion-induced
degradation~\citep{shumailov2023curse}---manifest precisely once an explicit critique loop is
added. The effect sizes are large and consistent: the single- vs.\ two-stage contrast
($d_z\approx-2.6$ to $-3.3$) dwarfs every within-two-stage manipulation.

\para{A caution about isolated LLM-judges.} The most striking observation is that absolute
quality scores rise slightly while meaning preservation falls sharply. This is a concrete
demonstration that an LLM-judge scoring documents \emph{in isolation} is blind to fidelity
drift---a caution for anyone using isolated LLM-judge scores as a convergence or quality
signal. Fidelity must be measured relative to the source, not judged from the current version
alone.

\para{Limitations.}
\begin{itemize}[leftmargin=*,itemsep=1pt,topsep=1pt]
    \item \textbf{Model/provider scope.} We use one reviser family (\gptone) with \gpto as
    cross-model reviewer and judge, due to key availability. The never-satisfied-reviewer
    mechanism is model-general in principle, but the specific non-convergence should be
    replicated on Claude, Gemini, and open-weight models before claiming universality.
    \item \textbf{Output-length cap.} The $2048$-token cap truncates later two-stage steps
    (mean $\sim 3$--$5$ of $8$ versions in the worst conditions). We mitigate with the
    truncation-free clean min $\Delta$, which still rejects convergence ($0.34\gg 0.02$,
    $p=5\times10^{-8}$); truncation is a \emph{downstream symptom} of inflation, not the cause
    of non-convergence. A larger cap would give cleaner late-step curves.
    \item \textbf{Horizon.} We cannot rule out convergence at much larger $T$, but the
    \texttt{D1\_main} $\dt$ curve is flat/rising from step $\sim 3$ and length is still growing
    at $V_8$, so there is no sign of approaching a fixed point.
    \item \textbf{Sample size.} $12$ documents per condition ($10$ for prose) is adequate for
    the very large primary effects (all $p<10^{-3}$) but modest for subtle ablation
    differences (harsh vs.\ structured is not significant).
    \item \textbf{Judge is an LLM.} Meaning preservation and quality are LLM-judged, not
    human-annotated; we reduce self-bias by using a different generation than the reviser, but
    human validation is future work.
    \item \textbf{Prompt dependence.} Results are for one reasonable \Review/\Revise prompt
    pair. A review prompt explicitly permitted to say ``no changes needed,'' or a reviser with
    a hard length budget, might converge---a positive-control follow-up.
\end{itemize}

\para{Broader implications.} Iterative review$\to$revise pipelines are deployed under the
implicit assumption that they stabilize at a polished document. Our results show the opposite:
the loop supplies an endless stream of edit pressure and degrades meaning while its own quality
signal reports success. Practitioners should not treat ``it stopped changing'' as an
achievable end state for these loops, and should not trust an isolated quality judge to detect
when to stop.
